\section{Introdução}

Incidentes cibernéticos que afetam hospitais e operadores de saúde têm impacto direto na continuidade do atendimento e na confidencialidade de dados de pacientes. Notícias especializadas registram esses eventos em linguagem natural e constituem uma fonte candidata para sistemas de extração de informações. Paim e Ignaczak \cite{paim2025} investigaram o uso de modelos pré-treinados \cite{devlin2019bert} para identificar entidades de incidentes cibernéticos em notícias da área da saúde e avaliaram SecureBERT-NER e SecureBERT-APTNER em uma amostra anotada manualmente. O corpus e as anotações originais, porém, não foram disponibilizados, o que impede a reprodução direta dos resultados.

Este artigo descreve a construção de um conjunto de dados próprio que preserva a tarefa e as quatro categorias do protocolo original, mantendo a distinção entre pré-anotação assistida e camada adjudicada. As contribuições são: (i)~um corpus de 299 notícias completas, documentado quanto à coleta; (ii)~duas camadas independentes, com 2.103 e 5.111 entidades; (iii)~avaliação sistemática de modelos e receitas de ajuste fino, com experimentos cruzados; e (iv)~uma ferramenta local de auditoria que permite categorizar entidades e adicionar trechos omitidos pela pré-anotação.

A seção~\ref{sec:corpus} descreve a construção do corpus e o protocolo de anotação. A seção~\ref{sec:exp} apresenta os experimentos e os resultados. A seção~\ref{sec:conc} discute limitações, disponibilidade e trabalhos futuros.
